% from ega2-indexes-and-contents-fr.tex
\oldpage[II]{215}
\markboth{TABLE DES MATIÈRES}{TABLE DES MATIÈRES}
\begin{longtable}{@{}p{0.09\textwidth}p{0.69\textwidth}r@{}}
\hyperref[subsection:II.8.7-fr]{8.7} & Éclatement des cônes projetants & 173\\
\hyperref[subsection:II.8.8-fr]{8.8} & Faisceaux amples et contractions & 177\\
\hyperref[subsection:II.8.9-fr]{8.9} & Le critère d'amplitude de Grauert : énoncé & 182\\
\hyperref[subsection:II.8.10-fr]{8.10} & Le critère d'amplitude de Grauert : démonstration & 184\\
\hyperref[subsection:II.8.11-fr]{8.11} & Unicité des contractions & 189\\
\hyperref[subsection:II.8.12-fr]{8.12} & Faisceaux quasi-cohérents sur les cônes projetants & 191\\
\hyperref[subsection:II.8.13-fr]{8.13} & Fermeture projective de sous-faisceaux et de sous-schémas fermés & 195\\
\hyperref[subsection:II.8.14-fr]{8.14} & Compléments sur les faisceaux associés aux $\mathcal{S}$-Modules gradués & 197\\
\multicolumn{2}{@{}l}{Bibliographie (suite)} & 205\\
\multicolumn{2}{@{}l}{Index des notations} & 207\\
\multicolumn{2}{@{}l}{Index terminologique} & 209\\
\multicolumn{2}{@{}l}{Errata et Addenda (Liste 1)} & 217\\
\end{longtable}

\begin{flushright}
\emph{Reçu le 22 juillet 1960.}
\end{flushright}
\endgroup
